\documentclass[
    a4paper,
    12pt,
    parskip=half
]{scrartcl}
\usepackage[utf8]{inputenc}
\usepackage[
    typ=pub,
    link=dlrg.net,
    contentLayout=bauchbinde,
    module={Paketbeschreibung}
]{dlrg}
\usepackage[
    language=ngerman,
    backend=biber,
    sortlocale=de_DE,%.UTF-8
    style=authoryear,
    bibencoding=UTF8,
    block=space,
    autocite=inline % \autocite[..][..]{..} erzeugt Literaturverweise mit runden Klammern
]{biblatex}
\usepackage{listings}
\usepackage{todonotes}
\usepackage{standalone}

\title{\LaTeX{}-Paket für die DLRG}
\subtitle{Abschnitt Publikation}
\author{Johannes Pieper}

\begin{document}
Der Dokumententyp \verbcode|pub| ist die Kurzform für Publikationen und zur Zeit die Standardeinstellung für Dokumntentypen des DLRG-Pakets. Dieser diese mehrseitigen Publikationen können z.\,B. Handbücher oder Regelwerke sein. Als Dokumentenklasse empfiehlt sich der Einsatz von \verbcode|scrartcl|.

\subsubsection{Layout}
Um das Layout dieser Publikation zu beeinflussen gibt es mehrere Optionen, die beim Paket mit angegeben werden können.
\begin{options}
    \keychoice{contentLayout}{bauchbinde,footline}\Default{footline}
        Über diese Option lässt sich das Gesamtlayout einstellen.
        Bei der ersten Möglichkeit bekommt jede Seite eine Bauchbinde. Die zweite Möglichkeit hat die einfache Wortmarke in der Fußzeile und ist gemäß CD-Handbuch \autocite[S.~13]{DLRG_CD2019} angelegt. Sie ist standardmäßig aktiviert.
    \keyval-{titelbild}{Dateiname}
        Wenn ein seitenfüllendes Titelbild eingebunden werden soll, so ist über diese Option der Dateiname anzugeben. Wird nichts angegeben erscheint der halbe Adler als Wasserzeichen auf der Titelseite
    \keyval{titelbildcaption}{Text}
        Ist ein Titelbild angegeben, so lässt sich hierrüber ein Bildtitel und damit auch ggf. verbunden der Verweis auf den Fotograph. Dieses wird dann in einem möglichen Abbildungsverzeichnis mit aufgeführt.
\end{options}

Zusätzlich kann die Titelseite nach eigenen Wünschen ergänzt werden. Auch kann das Format der Seite im Dokument gedreht werden.

\begin{commands}
    \command{titelExtra}[\marg{Titelelemente}] Ermöglicht auf der Titelseite weitere Elemente zu platzieren. Diese müssen als Parameter übergeben werden.
    \command{querformatSeite}
        Dreht das Standardmäßig als Hochformat eingestellt Format der Seite auf Querformat.
    \command{hochformatSeite}
        Dreht das Seitenformat wieder auf Hochformat zurück.
\end{commands}

\subsubsection{Boxen}
\begin{environments}
    \environment{redPartBox}
    Erzeugt eine rot umrandete Box, die Zwischenüberschriften enthalten kann.
    \begin{commands}
        \command{partTitle}[\marg{Überschrift}]
        Setzt die Zwischenüberschrift der \verbcode|redPartBox|.
    \end{commands}
    \begin{example}[gobble=8]
        \begin{redPartBox}
            \partTitle{Teil1}
            Etwas Text
            \partTitle{Teil2}
            Etwas mehr Text
        \end{redPartBox}
    \end{example}
\end{environments}
\end{document}
